% ISMRM abstract template: longitudinal T1 -> CT1 synthesis (I2SB + prior study), NYUMets.
% Structure follows the ISMRM submission form: Impact, Synopsis (Motivation / Goal(s) / Approach /
% Results), Abstract body (Introduction ... Conclusion), Acknowledgements, References, Figures.
% The form takes plain text pasted section by section; this file is for drafting, word counts and
% previewing the figures, not for upload.
%
% \todo{...} marks every number or decision still open. Figures are \input from figures/out/,
% written by figures/plot_main_bars.py --tikz, scripts/nyumets_volumes_csv.py --tikz,
% figures/table_main.py and figures/make_figure.py; rerun those and recompile.
% Limits (ISMRM, check the current call): Impact <= 40 words, Synopsis <= 100 words per field set,
% body <= 750 words, <= 5 figures. Word counts: texcount main.tex.

\documentclass[10pt]{article}
\usepackage[margin=2cm]{geometry}
\usepackage[T1]{fontenc}
\usepackage{lmodern}
\usepackage{amsmath,amssymb}
\usepackage{graphicx}
\usepackage{booktabs}
\usepackage{xcolor}
\usepackage{pgfplots}
\usepgfplotslibrary{groupplots}
\pgfplotsset{compat=1.16}
\usepackage{caption}
\captionsetup{font=small,labelfont=bf}
\usepackage[numbers,super,sort&compress]{natbib}

\newcommand{\todo}[1]{\textcolor{red}{[#1]}}
\newcommand{\figdir}{../../figures/out}
\newlength{\mainbarswidth}\setlength{\mainbarswidth}{\linewidth}
\newlength{\nyuprofilewidth}\setlength{\nyuprofilewidth}{0.7\linewidth}
% a framed placeholder until a figure file exists
\newcommand{\figorbox}[2]{\IfFileExists{#1}{\includegraphics[width=\linewidth]{#1}}%
  {\fbox{\parbox[c][4cm][c]{0.95\linewidth}{\centering\small\textcolor{red}{#2}}}}}

\begin{document}

% ------------------------------------------------------------------------------------------
\begin{center}
{\Large\bfseries \todo{Title}: Longitudinal Gadolinium-Free Contrast-Enhanced T1 Synthesis\\
from a Patient's Prior Exam with a Schr\"odinger Bridge}\\[6pt]
\todo{Eric \ldots}$^{1,2}$, \todo{co-authors}, Li Feng$^{1,2}$\\[2pt]
{\small $^1$\todo{Bernard and Irene Schwartz Center for Biomedical Imaging, Dept.\ of Radiology,
NYU Grossman School of Medicine, New York, NY, USA}\quad
$^2$\todo{CAI$^2$R, NYU Grossman School of Medicine, New York, NY, USA}}
\end{center}

% ------------------------------------------------------------------------------------------
\section*{Impact}
% <= 40 words
Brain-metastasis follow-up repeats contrast-enhanced MRI. Conditioning a contrast-synthesis
bridge on the patient's prior exam, registered without the scan being predicted, synthesizes
follow-up post-contrast T1 more accurately than the same network without it.
\todo{tumour-region result in one clause}

\section*{Synopsis}
\textbf{Motivation:} Patients under surveillance for brain metastases receive gadolinium at every
visit; their prior post-contrast exam is available but unused by contrast-synthesis methods.\\
\textbf{Goal(s):} Synthesize the current post-contrast T1 (CT1) from the current non-contrast
exam and the prior exam, using only information available at inference.\\
\textbf{Approach:} An image-to-image Schr\"odinger bridge (I2SB) maps T1 to CT1, conditioned on
FLAIR/T1/T2 and the prior T1/CT1 at five adjacent slices, re-registered from non-contrast
contrasts only; \todo{N} patients, 1007 test studies.\\
\textbf{Results:} The prior exam raised PSNR by 0.94\,dB and SSIM by 3.0 points ($p<10^{-100}$,
better in 97\% of studies). \todo{tumour region}

% ------------------------------------------------------------------------------------------
\section*{Introduction}
Brain metastases are monitored with serial gadolinium-enhanced MRI\cite{ref_rano}, and repeated
gadolinium exposure motivates synthesizing post-contrast images from non-contrast
ones\cite{ref_synth}. Single-exam synthesis must infer enhancement from T1, T2 and FLAIR alone.
In surveillance, however, every patient has a prior post-contrast exam, highly correlated with
the current one. We condition an image-to-image Schr\"odinger bridge (I2SB)\cite{ref_i2sb} on
that prior exam, and evaluate it under inference-realistic conditions: the prior is registered,
and all images are masked, without using the current CT1.

\section*{Methods}
\textbf{Data.} NYUMets\cite{ref_nyumets}: longitudinal brain-metastasis MRI (FLAIR, T1, CT1, T2)
\todo{scanner/protocol}. Studies of a patient were placed on a common grid by a 3-D affine fit
during preprocessing; each volume is scaled by its in-brain median, so prior and current share an
intensity scale. Training on slices 50--89 of \todo{N} patients; testing on slices 70--109 of
1007 studies from \todo{N} held-out patients with a prior exam.

\textbf{Model.} I2SB bridges the current T1 ($x_1$) to the current CT1 ($x_0$) with a
conditional U-Net (\todo{23M} parameters, effective batch 64, EMA, 300k steps), conditioned on
FLAIR/T1/T2. The longitudinal model additionally receives the prior T1 and CT1 at slices
$z\!-\!2\ldots z\!+\!2$ as input channels. Sampling: 20 steps.

\textbf{Inference-time registration and masking.} The preprocessing alignment saw the current
CT1 (through its brain and field-of-view masks), which inference cannot. Each prior is therefore
re-registered to the current study by a 3-D affine (rotation, per-axis scale, shift) maximising
the mean masked correlation of FLAIR, T1 and T2 of both exams, and all four prior contrasts are
resampled with it. Images and metrics are masked by the current T1 support intersected with the
registered prior CT1 support.

\textbf{Evaluation.} PSNR, SSIM and NRMSE within the mask, averaged per study; paired Wilcoxon
signed-rank tests over studies. \todo{Tumour-region error from the radiotherapy segmentations.}

\section*{Results}
Conditioning on the prior exam improved every metric (Figure~2, Table): PSNR
$24.79\!\pm\!0.09\rightarrow25.73\!\pm\!0.10$\,dB, SSIM $87.79\rightarrow90.83$\%, NRMSE
$13.07\rightarrow11.97$\% (all $p<10^{-160}$; better in 97\% of studies, median $+0.84$\,dB). The
gain grew from $+0.67$\,dB at the skull base (slice 70) to $\approx+1.05$\,dB above slice 95
(Figure~3). \todo{Tumour slices: gain inside lesions; new and resolved lesions (Figure~4).}
\todo{Example cases, Figure~1.}

\section*{Discussion}
\todo{Where the prior helps: stable anatomy vs.\ interval change; cases where it hurts.}
Limitations: \todo{the no-prior baseline was trained with fixed intensity scaling and on all
patients; one network family; slice-wise 2-D.}

\section*{Conclusion}
A patient's prior exam, registered without the scan being predicted, improves gadolinium-free
synthesis of the follow-up post-contrast T1. \todo{one clause on tumours}

\section*{Acknowledgements}
\todo{Grants.} The authors thank NYU Langone Health UltraViolet for technical support.

{\small
\begin{thebibliography}{9}
\bibitem{ref_rano} \todo{RANO-BM / surveillance imaging reference.}
\bibitem{ref_synth} \todo{Gadolinium-free contrast synthesis reference(s).}
\bibitem{ref_i2sb} Liu GH, Vahdat A, Huang DA, Theodorou EA, Nie W, Anandkumar A.
  I$^2$SB: Image-to-Image Schr\"odinger Bridge. In: Proc ICML. 2023. \todo{check}
\bibitem{ref_nyumets} \todo{NYUMets dataset reference.}
\end{thebibliography}}

% ------------------------------------------------------------------------------------------
\clearpage
\section*{Figures}

\begin{figure}[h]
\centering
\figorbox{\figdir/nyumets_longi/nyumets_longi.pdf}{\detokenize{figures/make_figure.py
figures/configs/nyumets_longi.py} (rows picked in the viewer)}
\caption{Example cases \todo{(rows)}: current T1 (input), registered prior CT1, synthesis without
and with the prior exam, and the acquired CT1; magnified region with its 5$\times$ residual,
PSNR/SSIM per panel.}
\end{figure}

\begin{figure}[h]
\centering
\input{\figdir/nyumets_bars.tikz}
\caption{1007 test studies, slices 70--109: mean $\pm$ s.e.\ over studies; *** paired Wilcoxon
$p\le0.001$. Value axes do not start at zero.}
\end{figure}

\begin{figure}[h]
\centering
\input{\figdir/nyumets_profile.tikz}\\[6pt]
{\small\input{\figdir/nyumets_table.tex}}
\caption{Top: PSNR per original slice index (mean $\pm$ s.e.\ over studies) and the paired gain
of the prior exam. Bottom: summary; bold best, * $p<0.05$ vs.\ the longitudinal model.}
\end{figure}

\begin{figure}[h]
\centering
\figorbox{\figdir/nyumets_tumor/nyumets_tumor.pdf}{Tumour cases: \detokenize{scripts/nyumets_tumor_slices.py
--write-rows}, then \detokenize{figures/make_figure.py}}
\caption{\todo{Lesion-containing slices: new, stable and resolved lesions; error inside the
lesion without and with the prior exam.}}
\end{figure}

\end{document}
